\documentclass[12pt,a4paper]{article}
\usepackage[margin=25mm, headheight=15pt, headsep=10pt]{geometry}
\usepackage{fontspec}
\usepackage[table]{xcolor} % Note: table option is required for rowcolors
\usepackage{titlesec}
\usepackage{tocloft}
\usepackage{fancyhdr}
\usepackage{booktabs}
\usepackage{tcolorbox}
\tcbuselibrary{skins, breakable}
\usepackage[colorlinks=true, linkcolor=emerald, urlcolor=emerald, bookmarks=true]{hyperref}

\definecolor{emerald}{RGB}{0,102,51}
\definecolor{darkred}{RGB}{139,0,0}
\definecolor{lightgray}{RGB}{245,245,245}

\usepackage[nil]{babel}
\babelprovide[import, main]{bengali}
\babelprovide[import]{arabic}
\babelfont{rm}[Renderer=HarfBuzz,Scale=1.0]{Noto Serif Bengali}
\babelfont{sf}[Renderer=HarfBuzz]{Noto Sans Bengali}
\babelfont{tt}[Renderer=HarfBuzz]{Noto Sans Bengali}
\babelfont[arabic]{rm}[Renderer=HarfBuzz,Scale=1.1]{Noto Naskh Arabic}

% Section styling
\titleformat{\section}{\Large\bfseries\color{emerald}}{\thesection.}{0.5em}{}
\titleformat{\subsection}{\large\bfseries\color{emerald!85!black}}{\thesubsection.}{0.5em}{}
\titleformat{\subsubsection}{\normalsize\bfseries\color{emerald!70!black}}{\thesubsubsection.}{0.5em}{}
\titlespacing*{\section}{0pt}{18pt}{8pt}

% Fancy Header/Footer setup
\pagestyle{fancy}
\fancyhf{} % clear all header and footer fields
\fancyhead[L]{\color{emerald}\small\textbf{\nouppercase{\leftmark}}}
\fancyfoot[L]{\scriptsize\color{black!50}{\lat{AI}}-সংকলিত, আলেম-যাচাইকৃত নয়}
\fancyfoot[C]{\color{emerald!70!black}\thepage}
\renewcommand{\headrulewidth}{0.4pt}
\renewcommand{\headrule}{\hbox to\headwidth{\color{emerald}\leaders\hrule height \headrulewidth\hfill}}
\renewcommand{\footrulewidth}{0pt}

% Custom boxes
% 1. Ayat/Hadith Box
\newtcolorbox{ayatbox}[1][]{
    enhanced, breakable, 
    colback=emerald!5, 
    colframe=emerald, 
    boxrule=0.5pt, 
    arc=3pt, 
    left=10pt, right=10pt, top=10pt, bottom=10pt,
    #1
}

% 2. Quote/Translation Box
\newtcolorbox{quotebox}[1][]{
    enhanced, breakable,
    frame hidden,
    colback=lightgray,
    borderline west={3pt}{0pt}{emerald},
    left=15pt, right=10pt, top=10pt, bottom=10pt,
    sharp corners,
    #1
}

% 3. AI-generated notice box
\newtcolorbox{warnbox}[1][]{
    enhanced, breakable,
    colback=red!4,
    colframe=darkred,
    boxrule=0.8pt,
    arc=3pt,
    left=10pt, right=10pt, top=10pt, bottom=10pt,
    #1
}

% Commands
\newcommand{\arab}[1]{\foreignlanguage{arabic}{#1}}
\newcommand{\kib}[1]{\textcolor{emerald}{\textbf{#1}}}

% Latin (English) fallback: Noto Serif Bengali has NO Latin glyphs
\newfontfamily\latfont[Renderer=HarfBuzz]{Noto Serif}
\newcommand{\lat}[1]{{\latfont #1}}

\setlength{\parskip}{5pt}
\setlength{\parindent}{0pt}

% Fix for missing bullet in Noto Serif Bengali
\renewcommand{\labelitemi}{$\bullet$}



\begin{document}
\begin{center}{\Large\bfseries\color{emerald} ঘটনা ৬ — উম্মে আতিয়্যাহ (রাঃ): ঈদে নারীদের জামাত-অংশগ্রহণ}\end{center}
\vspace{1em}
\section*{ঘটনা ৬ — উম্মে আতিয়্যাহ (রাঃ): ঈদে নারীদের জামাত-অংশগ্রহণ}

\begin{warnbox}
 \textbf{গুরুত্বপূর্ণ নোটিশ (\lat{Important} \lat{Notice}):} এই কন্টেন্টটি \lat{AI} (কৃত্রিম বুদ্ধিমত্তা) দিয়ে প্রস্তুত ও সংকলিত। \lat{AI} কঠোর সূত্র-মান নীতি মেনে শুধু নির্ভরযোগ্য উৎস থেকে তথ্য সংগ্রহ করেছে — কেবল সহীহ/হাসান হাদিস ও সহীহ সনদের আসার রাখা হয়েছে, দুর্বল সনদ বাদ দেওয়া হয়েছে। তবে এটি কোনো আলেম বা মুহাদ্দিস কর্তৃক যাচাইকৃত নয়।
\end{warnbox}

\begin{quote}
\textbf{সম্পর্কিত ফাইল:} পার্ট ৫গ (নারী-জামাত), পার্ট ৫ঘ (ঈদ)।
\end{quote}

\medskip\hrule\medskip

\subsection*{১. সূত্র কার্ড (\lat{Source} \lat{Card})}

\begin{table}[h]\centering\small
\begin{tabular}{ll}
বিষয় & বিবরণ \\
\hline
\textbf{ঘটনা} & উম্মে আতিয়্যাহ (রাঃ)-এর বর্ণনা — নবীজি (সা.) নারীদের (হায়েযা-সহ) ঈদের জামাত ও দোয়ায় অংশ নেওয়ার নির্দেশ দিয়েছেন \\
\textbf{রাবি} & উম্মে আতিয়্যাহ (রাঃ); হাফসা বিনতে সিরিন (রাবী হিসাবে) \\
\textbf{গ্রন্থ ও নম্বর} & সহীহ বুখারি ৯৭৪, ৯৮০, ৯৮১; সহীহ মুসলিম ৮৯০\lat{c} \\
\textbf{সনদের মান} & সহীহ \\
\hline
\end{tabular}
\end{table}

\medskip\hrule\medskip

\subsection*{২. পটভূমি (\lat{Background})}

\textbf{কে:} উম্মে আতিয়্যাহ আনসারিয়া (রাঃ) — নবীজির (সা.) যুগের সাহাবীয়া; অত্যন্ত সচেতন ও ধার্মিক নারী। হাফসা বিনতে সিরিন — তাবেয়িয়া।
\textbf{কখন:} মদিনা যুগ — বিশেষত রমজান-উত্তর ঈদুল ফিতর ও ঈদুল আজহার সময়; হাদিস-বর্ণনা দুই ঈদই উল্লেখ করে।
\textbf{কোথায়:} মদিনা — ঈদগাহ (ঈদের নামাজের স্থান)।
\textbf{কেন:} নবীজি (সা.) নারীদের ঈদ-গাহে অংশ নিয়ে আল্লাহর জিকির ও মুসলিম-দোয়ায় শরিক হতে উৎসাহিত করেছেন।

\medskip\hrule\medskip

\subsection*{৩. পূর্ণ ঘটনার বর্ণনা (\lat{The} \lat{Full} \lat{Event})}

উম্মে আতিয়্যাহ (রাঃ) বলেন:

\begin{quote}
\lat{“}নবীজি (সা.) আমাদেরকে আদেশ করেছিলেন — ঈদুল ফিতর ও ঈদুল আজহায় আমরা যেন\textbf{বালেগা ও কিশোরী (ওয়াতিক = যুবতী), পর্দানশীন ও হায়েযা} নারীদের বের করি। \textbf{হায়েযা নারীরা নামাজ (মসজিদ-প্রাঙ্গণ) থেকে আলাদা থাকবে — কিন্তু ঈদের জিকির-দোয়া ও মুসলিমদের সমাবেশে অংশ নেবে।}\lat{”}
\end{quote}

(অন্য বর্ণনায় — হাফসা বিনতে সিরিন থেকে)

নারীদের একজন প্রশ্ন করলেন:
\begin{quote}
\lat{“}হে আল্লাহর রাসূল! আমাদের কারো জিলবাব (পর্দার বাইরের কাপড়) না থাকলে — সে কি (ঈদে) যাবে না?\lat{”}
\end{quote}

নবীজি (সা.) বললেন:
\begin{quote}
\textbf{\lat{“}সে যেন তার বোন (সহচরী)-এর জিলবাব পরিধান করে — এবং ঈদের কল্যাণ ও মুসলিমদের দোয়ায় অংশ নেয়।\lat{”}}
\end{quote}

উম্মে আতিয়্যাহর (রাঃ) বর্ণনায় আরও আছে — হাফসা (রাঃ) জিজ্ঞেস করেছিলেন: \lat{“}(হায়েযা?)\lat{”} — উম্মে আতিয়্যাহ বললেন:
\begin{quote}
\lat{“}হ্যাঁ — হায়েযা নারী কি আরাফাত-এ (হজের ময়দানে) অংশ নেয় না?\lat{”} — অর্থাৎ হায়েয অবস্থায় নামাজ থেকে বিরত থাকা মানেই ইবাদত-সমাবেশ থেকে বের হওয়া নয় — ঈদের জিকির-দোয়ায় অংশ নেওয়া বৈধ ও নবীজির (সা.) নির্দেশ।
\end{quote}

\medskip\hrule\medskip

\subsection*{৪. ধাপে ধাপে বিশ্লেষণ (\lat{Step-by-Step} \lat{Analysis})}

\begin{enumerate}
\item \textbf{নারীদের ঈদ-অংশগ্রহণের নির্দেশ:} নবীজি (সা.) নারীদের সমাবেশ থেকে আলাদা থাকতে নিষেধ করেননি — বরং ঈদের জিকির-দোয়ায় শরিক হওয়ার আদেশ দিয়েছেন।
\item \textbf{হায়েযা-নারীর নিয়ম:} হায়েযা নারী নামাজ (সালাতে) অংশ নেবে না — কারণ তা বৈধ নয় (পার্ট ৫গ) — কিন্তু ঈদের সমাবেশ-জিকির-দোয়ায় অংশ নেবে — এ দুয়ের পার্থক্যই নবীজির (সা.) শিক্ষা।
\item \textbf{জিলবাব-অভাবের সমাধান:} সামাজিক সংহতি — একজনের না থাকলে অন্য সহচরীর দিয়ে পূরণ — সম্পদের অভাব বাধা নয়, ব্যবস্থা শিক্ষা।
\item \textbf{নবী-যুগের নারীর ঈদ-ইবাদত:} মসজিদে নববি-নির্ভর নয় — ঈদগাহ-সহ নবীজির যুগের নারীদের উৎসাহের প্রামাণ্য — নারীদের মসজিদ-জামাতের বৈধতার ভিত্তি (পার্ট ৫ক)।
\end{enumerate}
\medskip\hrule\medskip

\subsection*{৫. শব্দের ব্যাখ্যা (\lat{Key} \lat{Words})}

\begin{itemize}
\item \textbf{ওয়াতিক (ʿ\lat{Awatiq}):} যুবতী-বালিগা নারী — ঈদে বের হওয়ার যোগ্য।
\item \textbf{যাওয়াতুল খুদুর:} পর্দায়-অন্তরালে থাকা নারী — ঘরের ভেতরে-পর্দার আড়ালে থাকা।
\item \textbf{হায়েয:} মাসিক — যাতে নামাজ-রোজা-কুরআন-স্পর্শের (মাযহাবগত) বিধান প্রযোজ্য, কিন্তু জিকির-দোয়া থেকে বিরত নয়।
\item \textbf{জিলবাব:} বাইরের পোশাক (চাদর) — সতর ঢাকার জন্য।
\end{itemize}
\medskip\hrule\medskip

\subsection*{৬. সংশ্লিষ্ট হাদিস ও আয়াত}

\begin{quote}
\textbf{\lat{“}হে নবী! আপনার স্ত্রী, কন্যাদের ও মুমিন নারীদের বলুন — তারা যেন তাদের চাদরের (জিলবাব) অংশ নিজেদের উপর টেনে দেয়। এতে তাদের চেনা সহজতর হবে — ফলে তাদেরকে কষ্ট দেওয়া হবে না।\lat{”}} — সূরা আল-আহযাব ৩৩:৫৯ (বাংলা অর্থ) — জিলবাব-প্রসঙ্গে নবীজির উত্তর।
\end{quote}

\begin{quote}
\textbf{\lat{“}আর মুমিন নারীদের বলুন — তারা যেন তাদের দৃষ্টি সংযত রাখে, লজ্জাস্থানের হেফাজত করে...\lat{”}} — সূরা আন-নূর ২৪:৩১ (বাংলা অর্থ)।
\end{quote}

\medskip\hrule\medskip

\subsection*{৭. গভীর শিক্ষা (\lat{Deep} \lat{Lessons})}

\textbf{শিক্ষা ১ — নারীর ইবাদত-অংশগ্রহণ:} নারীদের নামাজ-মসজিদ থেকে আলাদা রাখা হলো এক বিষয়; কিন্তু ঈদ-জিকির, দোয়া, জ্ঞান-চর্চা (ইলম) — সবসময় উৎসাহিত — নবীযুগের প্রমাণ।

\textbf{শিক্ষা ২ — শর্ত-মেনে অংশগ্রহণ:} হায়েযায় নারী নামাজ-স্থানে নয় — কিন্তু সমাবেশে; অর্থাৎ \textbf{বিচ্ছিন্নতা নয়, সীমিততা} — ইসলাম নারীকে ঘর-বন্দি শিক্ষা দেয়নি।

\textbf{শিক্ষা ৩ — সামাজিক সমর্থন:} জিলবাব-অভাবের শিশু-নারীও বাদ পড়েনি — পারস্পরিক সহযোগিতা (সহচরীর জিলবাব) — সমাজের সংহতির শিক্ষা।

\textbf{শিক্ষা ৪ — ঈদের জন্য পোশাক-প্রস্তুতি:} সতর-ঢাকা বা উত্তম পোশাক — ঈদের পরিচ্ছন্নতা-সৌন্দর্য — নবীজির (সা.) শিক্ষা।

\medskip\hrule\medskip

\subsection*{৮. জীবনে প্রয়োগের পদ্ধতি (\lat{Practical} \lat{Application})}

\begin{enumerate}
\item \textbf{ঈদের আমল:} নারীরাও সতর-ঢাকা, পরিচ্ছন্ন পোশাকে ঈদ-জামাত ও ঈদগাহ-সমাবেশে অংশ নিতে পারেন (যথাসম্ভব পর্দা-রক্ষায় আলাদা স্থান)।
\item \textbf{হায়েযা-নারীদের জন্য:} নামাজ থেকে বিরত থাকবেন কিন্তু ঈদ-দোয়া, জিকির, ইসলামিক জ্ঞান-সভায় (মহিলা সমাবেশ) অংশ নেওয়া — নবীজির (সা.) নির্দেশ মান্য।
\item \textbf{সহযোগিতার মনোভাব:} কারো পোশাক-অভাব থাকলে তাকে সাহায্য করা (জিলবাব-ধার) — ঈদ-সমাবেশে কাউকে বাদ না দেওয়া।
\item \textbf{ঈদের প্রস্তুতি:} ঈদের সকালে পরিচ্ছন্নতা, সেরা পোশাক, তাকবির-গণ — নারী-পুরুষ উভয়ে।
\end{enumerate}
\medskip\hrule\medskip

\subsection*{৯. রেফারেন্স}

\begin{itemize}
\item সহীহ বুখারি — ৯৭৪, ৯৮০, ৯৮১
\item সহীহ মুসলিম — ৮৯০\lat{c}
\item সুনানে আবু দাউদ — ১১৫৪-এর মতো ঈদ-সংক্রান্ত
\item কুরআন: আল-আহযাব ৩৩:৫৯; আন-নূর ২৪:৩১
\end{itemize}

\end{document}
